LLM pricing agents may soon publish a chart beside each price, while audits of algorithmic pricing read only declared prices and text. Antitrust has met this gap before: in the ATPCO case, airlines were alleged to signal through fare-filing footnotes that carried no fare data. We test whether an LLM seller can pass a target price to a rival through a chart whose data are locked to its declared prices, whether monitors notice, and what removes the signal. Told to signal one of 11 target prices with no agreed code, gpt-6.1-sol drew charts that it decoded in \NschemelessSolSol\ trials and that gpt-5.5 decoded in \NextSolFive\ (chance 1/11). Four prompted LLM monitors flagged \Nsolncmin--\Nsolncmax\% of these charts at a 5\% false-positive rate. Detectors that know the target grid or the mark do better, but in a matched control of honest charts with a decorative element the grid-aware decoder raised a false alarm on \NdecorFD, and trained detectors could not tell these charts from signalling ones; only a rule written for the observed mark stayed selective. Redrawing every chart from the declared data left \Nredrawall\ correct decodes, and a one-line proof shows that no decoder can beat chance on a redraw. In a \Ngamerounds-round pilot game with no instruction to signal, sellers imitated each other's raw chart styling but did not detectably raise prices. The remedy used for fare filings, publishing only from declared data, also works for agent charts.
